\documentclass[11pt]{article}
\usepackage[T1]{fontenc}
\usepackage[utf8]{inputenc}
\usepackage{lmodern}
\usepackage[letterpaper,margin=0.78in]{geometry}
\usepackage{microtype}
\usepackage{booktabs,longtable,tabularx,array,multirow}
\usepackage{amsmath,amssymb}
\usepackage{xcolor}
\usepackage{enumitem}
\usepackage{titlesec}
\usepackage{fancyhdr}
\usepackage{lastpage}
\usepackage{hyperref}
\usepackage[most]{tcolorbox}
\usepackage{parskip}
\usepackage{needspace}
\usepackage{ragged2e}
\usepackage{setspace}

\definecolor{Navy}{HTML}{17365D}
\definecolor{Blue}{HTML}{2F75B5}
\definecolor{Teal}{HTML}{1F6D6A}
\definecolor{Gold}{HTML}{B38B2D}
\definecolor{Gray}{HTML}{5E6B75}
\definecolor{LightBlue}{HTML}{EEF4FA}
\definecolor{LightGold}{HTML}{FBF6E8}
\definecolor{LightGray}{HTML}{F4F5F7}
\definecolor{LightTeal}{HTML}{EDF7F6}

\hypersetup{colorlinks=true,linkcolor=Navy,urlcolor=Blue,citecolor=Teal,pdfauthor={OpenAI},pdfsubject={Value investing study and implementation materials}}
\setlist[itemize]{leftmargin=1.5em,itemsep=2pt,topsep=3pt}
\setlist[enumerate]{leftmargin=1.7em,itemsep=2pt,topsep=3pt}
\renewcommand{\arraystretch}{1.16}
\setlength{\tabcolsep}{5pt}
\titleformat{\section}{\Large\bfseries\color{Navy}}{\thesection}{0.6em}{}[\vspace{-0.25em}\color{Blue}\rule{\linewidth}{0.5pt}]
\titleformat{\subsection}{\large\bfseries\color{Teal}}{\thesubsection}{0.5em}{}
\titleformat{\subsubsection}{\normalsize\bfseries\color{Gray}}{\thesubsubsection}{0.5em}{}
\pagestyle{fancy}
\fancyhf{}
\fancyhead[L]{\small\color{Gray}\DocShortTitle}
\fancyhead[R]{\small\color{Gray}Value Investing Reference Set}
\fancyfoot[C]{\small\color{Gray}Page \thepage\ of \pageref*{LastPage}}
\renewcommand{\headrulewidth}{0.4pt}
\renewcommand{\footrulewidth}{0pt}

\newtcolorbox{principlebox}[1][]{enhanced,breakable,colback=LightBlue,colframe=Blue,boxrule=0.6pt,arc=2mm,left=2mm,right=2mm,top=1.5mm,bottom=1.5mm,title={#1},fonttitle=\bfseries}
\newtcolorbox{cautionbox}[1][]{enhanced,breakable,colback=LightGold,colframe=Gold,boxrule=0.6pt,arc=2mm,left=2mm,right=2mm,top=1.5mm,bottom=1.5mm,title={#1},fonttitle=\bfseries}
\newtcolorbox{modernbox}[1][]{enhanced,breakable,colback=LightTeal,colframe=Teal,boxrule=0.6pt,arc=2mm,left=2mm,right=2mm,top=1.5mm,bottom=1.5mm,title={#1},fonttitle=\bfseries}
\newtcolorbox{worksheetbox}[1][]{enhanced,breakable,colback=white,colframe=Gray,boxrule=0.45pt,arc=1mm,left=2mm,right=2mm,top=1.5mm,bottom=1.5mm,title={#1},fonttitle=\bfseries}
\newcommand{\checkbox}{\(\square\)}
\newcommand{\blankline}{\rule{0.96\linewidth}{0.35pt}}
\newcommand{\sourceRef}[1]{\textit{Source anchor: #1}}
\newcolumntype{Y}{>{\RaggedRight\arraybackslash}X}
\newcolumntype{C}[1]{>{\centering\arraybackslash}p{#1}}

\newcommand{\makedoccover}[4]{%
\begin{titlepage}
\thispagestyle{empty}
\vspace*{1.0in}
{\color{Navy}\rule{\linewidth}{2.2pt}}\\[0.55in]
{\Huge\bfseries\color{Navy} #1\par}
\vspace{0.18in}
{\Large\color{Teal} #2\par}
\vspace{0.5in}
\begin{tcolorbox}[colback=LightGray,colframe=Gray,boxrule=0.5pt,arc=2mm]
\textbf{Purpose.} #3
\end{tcolorbox}
\vfill
{\large Prepared from the revised edition of \textit{The Intelligent Investor}, with Jason Zweig's commentary.\par}
\vspace{0.12in}
{\small #4\par}
\vspace{0.25in}
{\small\color{Gray}\parbox{0.90\textwidth}{\centering Educational material only. Historical thresholds and examples are preserved for study; they are not individualized investment, tax, or legal advice.}}\par
\vspace{0.35in}
{\color{Navy}\rule{\linewidth}{2.2pt}}
\end{titlepage}
}

\newcommand{\DocShortTitle}{Framework \& Formulas}
\begin{document}
\makedoccover{Graham Investment Framework and\\Formula Reference}{Screening, valuation, margin-of-safety, and portfolio arithmetic}{Separates source formulas, algebraic derivations, and modern implementation formulas.}{Standalone LaTeX study document}
\tableofcontents\newpage

\begin{cautionbox}[Source and modernization boundary]
The underlying source is the fourth revised edition, with Benjamin Graham's text updated in 1971--1972 and later commentary by Jason Zweig. Historical yields, tax rules, company examples, market levels, and dollar thresholds are therefore period-specific. This document preserves the source's framework and labels later implementation ideas separately rather than silently treating historical figures as current facts.
\end{cautionbox}

\section{Formula Map}
\begin{center}
\small
\begin{tabularx}{\linewidth}{@{}p{0.25\linewidth}p{0.30\linewidth}Y@{}}
\toprule
\textbf{Concept} & \textbf{Expression} & \textbf{Interpretation}\\
\midrule
Current ratio & $CR=CA/CL$ & Liquidity buffer. Chapter 14's industrial-company defensive threshold is $CR\ge 2$.\\
Net current assets & $NCA=CA-CL$ & Working-capital buffer before long-term claims.\\
Net-current-asset value & $NCAV=CA-TL-$ prior claims & Asset-based bargain measure; fixed assets receive no value in the stripped-down test.\\
Earnings yield & $EY=EPS/P=1/(P/E)$ & Earning power per dollar of purchase price.\\
Dividend yield & $DY=D/P$ & Cash distribution per dollar of market price.\\
Combined value screen & $(P/E)(P/B)\le22.5$ & Chapter 14 tradeoff between earnings and asset valuation.\\
Derived Graham Number & $P_{max}=\sqrt{22.5\,EPS\,BVPS}$ & Algebraic shorthand derived from the combined Chapter 14 thresholds; not separately named in the source.\\
Growth approximation & $V=E(8.5+2g)$ & Historical Chapter 11 growth-stock approximation; explicitly not a true-value formula.\\
Margin-of-safety spread & $MOS_s=EY-Y_b$ & One Chapter 20 expression of stock earning power versus bond yield.\\
DCA shares & $q_t=C_t/P_t$ & Shares acquired by periodic contribution.\\
Average cost & $\bar P=\sum C_t/\sum q_t$ & Effective average purchase price under periodic investing.\\
Rebalancing trade & $\Delta_i=w_i^*V-w_iV$ & Dollar amount needed to restore target weight $w_i^*$.\\
\bottomrule
\end{tabularx}
\end{center}

\section{Defensive Financial-Strength Tests}
\subsection{Industrial liquidity}
\begin{align}
CR &= \frac{Current\ Assets}{Current\ Liabilities}\ge 2.0,\\
LongTermDebt &\le Current\ Assets-Current\ Liabilities.
\end{align}
These are source-era tests. The modern analytical principle is that liquidity and long-term leverage should both leave room for an adverse operating period.

\subsection{Earnings and dividend record}
The defensive screen asks for positive earnings in every one of ten years and a long uninterrupted dividend history. The source also requires at least a one-third increase in per-share earnings over a decade using three-year averages at the starting and ending points.

A useful way to express the growth test is:
\begin{equation}
Growth_{10}=\frac{\overline{EPS}_{end,3}}{\overline{EPS}_{start,3}}-1 \ge \frac{1}{3}.
\end{equation}
A corresponding annualized rate can be computed as:
\begin{equation}
CAGR=\left(\frac{\overline{EPS}_{end,3}}{\overline{EPS}_{start,3}}\right)^{1/n}-1,
\end{equation}
where $n$ is the year distance between the midpoints of the two averaging windows. This CAGR expression is a modern analytical convenience; the source's criterion itself is the one-third total increase.

\section{Valuation Tests}
\subsection{P/E limit}
\begin{equation}
P \le 15\times \overline{EPS}_{3yr}.
\end{equation}
This uses normalized multi-year earnings rather than one favorable year.

\subsection{Price/book limit}
\begin{equation}
P \le 1.5\times BVPS.
\end{equation}
The source permits tradeoff between P/E and P/B so long as the product does not exceed 22.5.

\subsection{Graham Number derivation}
Starting from
\[
\frac{P}{EPS}\cdot\frac{P}{BVPS}\le22.5,
\]
then
\[
P^2\le22.5\,EPS\,BVPS,
\]
which gives
\[
P\le\sqrt{22.5\,EPS\,BVPS}.
\]
\begin{cautionbox}[Interpretation limit]
This shortcut is sensible only when both earnings and book value are economically meaningful. It can be misleading for financial firms with unusual accounting, asset-light firms whose durable economics reside mainly in intangible capital, cyclical firms at peak earnings, or firms with impaired assets.
\end{cautionbox}

\section{Growth-Stock Approximation}
The Chapter 11 expression is:
\begin{equation}
V=E(8.5+2g).
\end{equation}
\textbf{Variables:} $E$ is current normal earnings per share; $g$ is the expected annual growth rate over approximately seven to ten years, expressed as a whole-number percentage in the historical presentation (e.g., 5 rather than 0.05).

\begin{cautionbox}[Do not treat as intrinsic-value truth]
The book's own endnote says the expression is not offered as the true value of a growth stock. Long-range growth estimates dominate the output, so small forecast errors can cause large valuation errors.
\end{cautionbox}

\section{Bargain and NCAV Tests}
\subsection{General bargain threshold}
Chapter 7 suggests that an issue should not be called a true bargain unless indicated value is at least 50\% above market price:
\begin{equation}
V_{indicated} \ge 1.5P.
\end{equation}
Equivalently, price is no more than two-thirds of indicated value.

\subsection{Net-current-asset method}
\begin{align}
NCAV &= Current\ Assets-Total\ Liabilities-Prior\ Claims,\\
P_{portfolio} &\lesssim \frac{2}{3}NCAV \quad \text{(historical Graham-Newman practice)}.
\end{align}
The method depends on diversification because liquidation values, accounting quality, and business deterioration can differ materially from the screen.

\section{Margin of Safety}
\subsection{Bond coverage form}
For a fixed claim, define:
\begin{equation}
Coverage=\frac{Earnings\ Available\ for\ Fixed\ Charges}{Fixed\ Charges}.
\end{equation}
The safety margin is the amount by which earnings can decline before coverage falls to 1.0. A simple percentage form is:
\begin{equation}
MOS_{coverage}=1-\frac{1}{Coverage}.
\end{equation}
This percentage is a derived diagnostic, not a quoted Graham formula.

\subsection{Equity earning-power spread}
\begin{equation}
MOS_s=\frac{EPS}{P}-Y_b.
\end{equation}
The source treats the excess of stock earning power over bond interest as one form of cushion. This is not sufficient by itself because future earnings can decline and accounting earnings can overstate distributable economics.

\subsection{Valuation-discount form}
A modern generic expression is:
\begin{equation}
MOS_v=\frac{V-P}{V}.
\end{equation}
This is a useful reporting convention for a value estimate $V$, but it is not a substitute for conservatism in how $V$ was estimated.

\section{Portfolio Formulas}
\subsection{Allocation and rebalance}
For a target weight $w_i^*$ and portfolio value $V$:
\begin{equation}
TargetDollar_i=w_i^*V, \qquad Trade_i=TargetDollar_i-CurrentDollar_i.
\end{equation}
A positive trade means buy; negative means sell. Taxable-account implementation should consider tax consequences before mechanical selling.

\subsection{Contribution-directed rebalancing}
When possible, use new cash to reduce underweights:
\begin{equation}
AllocationOfNewCash_i \propto \max(0,w_i^*-w_i).
\end{equation}
This modern implementation minimizes turnover while preserving the source's policy-discipline principle.

\section{Risk and Recovery Arithmetic}
Losses are asymmetric with required recovery gains:
\begin{equation}
RequiredGain=\frac{L}{1-L},
\end{equation}
where $L$ is the fractional loss. A 50\% loss requires a 100\% gain; a 75\% loss requires a 300\% gain; a 95\% loss requires a 1,900\% gain. This arithmetic explains why avoiding catastrophic permanent loss is more important than maximizing upside in every period.

\section{Formula Governance}
Before using any metric, document:
\begin{enumerate}
\item the accounting period and whether earnings are normalized;
\item whether book value is tangible and economically relevant;
\item whether debt-like obligations are included consistently;
\item whether the metric is a source rule, a derived identity, or a modern implementation;
\item which conditions invalidate the metric;
\item what independent qualitative evidence is required before capital is committed.
\end{enumerate}

\section*{Primary Source}
Benjamin Graham, \textit{The Intelligent Investor: A Book of Practical Counsel}, revised edition, updated with commentary by Jason Zweig, preface and appendix material by Warren E. Buffett. Chapter references in this document point to the structure of that edition.

\end{document}
